% sections/sorting.tex
\chapter{Алгоритм сортировки слиянием}

\section{Теоретическая часть}

Сортировка слиянием (\textit{Merge Sort}) — рекурсивный алгоритм сортировки,
предложенный Джоном фон Нейманом в~1945 году~\cite{knuth1997}. Алгоритм
разбивает массив на два подмассива, рекурсивно сортирует их, а затем
сливает в один отсортированный массив.

\subsection{Вычислительная сложность}

Пусть \(T(n)\)~--- время работы алгоритма на массиве длины~\(n\).
Рекуррентное соотношение имеет вид:

\begin{equation}
  \label{eq:merge-recurrence}
  T(n) = 2\,T\!\left(\frac{n}{2}\right) + \Theta(n),
  \quad T(1) = \Theta(1).
\end{equation}

По основной теореме о рекуррентных соотношениях (теорема Мастера)
решением~\eqref{eq:merge-recurrence} является:

\begin{equation}
  \label{eq:merge-complexity}
  T(n) = \Theta(n \log n).
\end{equation}

Таким образом, сортировка слиянием работает за~\(\Theta(n \log n)\) во всех
случаях, что является асимптотически оптимальным для алгоритмов, основанных
на сравнении~\cite{cormen2022}.

Количество сравнений при слиянии двух подмассивов длин \(p\) и~\(q\)
не превышает:

\begin{equation}
  \label{eq:merge-comparisons}
  C(p, q) \leq p + q - 1.
\end{equation}

\subsection{Сравнение алгоритмов сортировки}

В таблице~\ref{tab:sort-comparison} приведено сравнение временной и
пространственной сложности нескольких алгоритмов сортировки~\cite{sedgewick2011}.

\begin{table}[H]
  \centering
  \caption{Сравнение алгоритмов сортировки}
  \label{tab:sort-comparison}
  \begin{tabular}{lcccc}
    \toprule
    \textbf{Алгоритм}
      & \textbf{Лучший случай}
      & \textbf{Средний случай}
      & \textbf{Худший случай}
      & \textbf{Память} \\
    \midrule
    Пузырьковая     & $\Omega(n)$       & $\Theta(n^2)$       & $O(n^2)$       & $O(1)$ \\
    Быстрая         & $\Omega(n\log n)$ & $\Theta(n\log n)$   & $O(n^2)$       & $O(\log n)$ \\
    \textbf{Слиянием} & $\Omega(n\log n)$ & $\Theta(n\log n)$ & $O(n\log n)$   & $O(n)$ \\
    Пирамидальная   & $\Omega(n\log n)$ & $\Theta(n\log n)$   & $O(n\log n)$   & $O(1)$ \\
    Вставками       & $\Omega(n)$       & $\Theta(n^2)$       & $O(n^2)$       & $O(1)$ \\
    \bottomrule
  \end{tabular}
\end{table}

\subsection{Структура рекурсии}

На рисунке~\ref{fig:merge-tree} изображено дерево рекурсивных вызовов
сортировки слиянием для массива из семи элементов.

\begin{figure}[H]
  \centering
  \includegraphics[width=0.92\textwidth]{images/merge_sort_tree.png}
  \caption{Дерево рекурсии алгоритма сортировки слиянием}
  \label{fig:merge-tree}
\end{figure}

\section{Реализация на Python}

Ниже представлена реализация сортировки слиянием на языке Python~\cite{python_docs}.
Функция \mintinline{python}{merge_sort} принимает список \mintinline{python}{arr}
и возвращает новый отсортированный список. Вспомогательная функция
\mintinline{python}{merge} сливает два упорядоченных списка
(\mintinline{python}{left} и \mintinline{python}{right}) в один.

\inputminted[
  label=\fbox{merge\_sort.py},
  firstline=1,
  lastline=50
]{python}{code/merge_sort.py}

\noindent
Листинг~\ref{code:merge-sort} расположен в приложении~\ref{app:merge-sort-full}
(полный исходный файл).
